% TODO make hw stylesheet
\documentclass[10pt]{article}
\usepackage[letterpaper,top=1in]{geometry}
\usepackage[dvipsnames,svgnames]{xcolor}
\usepackage[shortlabels]{enumitem}
\usepackage[framemethod=TikZ]{mdframed}
\usepackage{amsmath,amssymb,amsthm}
\usepackage[colorlinks]{hyperref}
\usepackage{multicol}
\usepackage{tabularx}

\hypersetup{
	colorlinks=true,
	linkcolor=blue,
	filecolor=magenta,
	urlcolor=magenta,
	pdftitle={MAT 528 HW}
}

\usepackage{mathtools}
\usepackage{thmtools}
\usepackage[textsize=footnotesize]{todonotes}
\usepackage{titling}

\reversemarginpar
\mdfdefinestyle{mdbluebox}{roundcorner=10pt,innerbottommargin=9pt,
	linecolor=blue,backgroundcolor=TealBlue!5,}
\declaretheoremstyle[headfont=\sffamily\bfseries\color{MidnightBlue},
mdframed={style=mdbluebox},]{thmbluebox}

\mdfdefinestyle{mdbloobox}{frametitlefont=\bfseries,innerbottommargin=8pt,
	nobreak=true,backgroundcolor=TealBlue!5,linecolor=blue,}
\declaretheoremstyle[headfont=\bfseries\color{MidnightBlue},
mdframed={style=mdbloobox},headpunct={\\[3pt]},postheadspace=0pt,]{thmbloobox}

\mdfdefinestyle{mdredbox}{frametitlefont=\bfseries,innerbottommargin=8pt,
	nobreak=true,backgroundcolor=Salmon!5,linecolor=RawSienna,}
\declaretheoremstyle[headfont=\bfseries\color{RawSienna},
mdframed={style=mdredbox},headpunct={\\[3pt]},postheadspace=0pt,]{thmredbox}

\mdfdefinestyle{mdgreenbox}{linecolor=ForestGreen,backgroundcolor=ForestGreen!5,
	linewidth=2pt,rightline=false,leftline=true,topline=false,bottomline=false,}
\declaretheoremstyle[headfont=\bfseries\sffamily\color{ForestGreen!70!black},
mdframed={style=mdgreenbox},headpunct={ --- },]{thmgreenbox}

\mdfdefinestyle{mdorangebox}{linecolor=RedOrange,backgroundcolor=RedOrange!5,
	linewidth=2pt,rightline=false,leftline=true,topline=false,bottomline=false,}
\declaretheoremstyle[headfont=\bfseries\sffamily\color{RedOrange!70!black},
mdframed={style=mdorangebox},headpunct={ --- },]{thmorangebox}

\mdfdefinestyle{mdblackbox}{linecolor=black,backgroundcolor=RedViolet!5!gray!5,
	linewidth=3pt,rightline=false,leftline=true,topline=false,bottomline=false,}
\declaretheoremstyle[mdframed={style=mdblackbox}]{thmblackbox}

\declaretheoremstyle[spaceabove=6pt,spacebelow=6pt]{basehead}

\declaretheorem[style=basehead,name=Problem]{problem}
\declaretheorem[style=thmredbox,name=Problem,numbered=no]{problem*}
\declaretheorem[style=thmbloobox,name=Setup]{setup} % for config geo
\declaretheorem[style=thmredbox,name=Example,sibling=problem]{example}
\declaretheorem[style=thmredbox,name=$\bigstar$ Problem,sibling=problem]{niceprob}
\declaretheorem[style=thmbluebox,name=Theorem,sibling=problem]{theorem}
\declaretheorem[style=thmbluebox,name=Corollary,sibling=theorem]{corollary}
\declaretheorem[style=thmbluebox,name=Proposition,sibling=theorem]{prop}
\declaretheorem[style=thmbluebox,name=Lemma,numberwithin=problem]{lemma}
\declaretheorem[style=thmbluebox,name=Theorem,numbered=no]{theorem*}
\declaretheorem[style=thmbluebox,name=Lemma,numbered=no]{lemma*}
\declaretheorem[style=thmgreenbox,name=Claim,numberwithin=problem]{claim}
\declaretheorem[style=thmgreenbox,name=Claim,numbered=no]{claim*}
\declaretheorem[style=thmblackbox,name=Remark,numberwithin=problem]{remark}
\declaretheorem[style=thmblackbox,name=Remark,numbered=no]{remark*}
\declaretheorem[style=thmgreenbox,name=Definition,sibling=theorem]{definition}
\declaretheorem[style=thmgreenbox,name=Definition,numbered=no]{definition*}

\providecommand{\ol}{\overline}
\providecommand{\ceil}[1]{\left\lceil #1 \right\rceil}
\providecommand{\floor}[1]{\left\lfloor #1 \right\rfloor}
\newcommand{\dd}{\mathrm{d}}
\providecommand{\ul}{\underline}
\providecommand{\eps}{\varepsilon}
\providecommand{\half}{\frac{1}{2}}
\providecommand{\CC}{\mathbb C}
\providecommand{\FF}{\mathbb F}
\providecommand{\NN}{\mathbb N}
\providecommand{\QQ}{\mathbb Q}
\providecommand{\RR}{\mathbb R}
\providecommand{\ZZ}{\mathbb Z}
\providecommand{\dg}{^\circ}
\providecommand{\ii}{\item}
\providecommand{\setdiff}{\smallsetminus}
\providecommand{\alert}[1]{{\sffamily\textbf{\textcolor{blue}{#1}}}}
\providecommand{\inv}{^{-1}}
\providecommand{\mb}[1]{\mathbf{#1}}

\newenvironment{soln}{\begin{proof}[Solution]}{\end{proof}}

\providecommand{\pow}{\operatorname{Pow}}
\providecommand{\vocab}[1]{{\textbf{\textcolor{ForestGreen}{#1}}}}
\providecommand{\lcm}{\operatorname{lcm}}
\providecommand{\abs}[1]{\left\lvert #1\right\rvert}
\providecommand{\norm}[1]{\left\|#1\right\|}

\usepackage{mathrsfs}

% place title at top right
\setlength{\droptitle}{-8ex}
\pretitle{\begin{flushright}\Large\bfseries}
\posttitle{\par\end{flushright}}
\preauthor{\begin{flushright}}
\postauthor{\end{flushright}}
\predate{\begin{flushright}}
\postdate{\end{flushright}}

\title{MAT 528 HW11}
\author{\large Aditya Pahuja}
\date{\large Due 12/01}
\begin{document}
\maketitle

\begin{problem}
    [Munkres 53.3]
    Let $p\colon E\to B$ be a covering map; let $B$ be connected.
    Show that if $p\inv(b_0)$ has $k$ elements for some $b_0\in B$,
    then $p\inv(b)$ has $k$ elements for every $b\in B$.
\end{problem}

\begin{soln}
    Given $b\in B$, let $U\subseteq B$ be a neighborhood of $b$
    that is evenly covered by $p$. The slices of $p\inv(U)$
    are in bijection with $p\inv(b)$ by definition,
    so, for each $b'\in U$, $p\inv(b')$ has the same cardinality
    as $p\inv(b)$ (with bijection between the two
    given by sending $x\in p\inv(b)$ to the $y\in p\inv(b')$
    that is in the same sheet of $p\inv(U)$ as $x$).
    This means that the cardinality of $p\inv(b)$
    is a locally constant function of $b$.
    Locally constant functions on connected spaces are constant,
    so the cardinality of $p\inv(b)$ is equal to
    that of $p\inv(b_0)$ for all $b\in B$,
    thus $p\inv(b)$ has $k$ elements, too.
\end{soln}

\begin{problem}
    [Munkres 53.5]
    Show that the map $p\colon S^1\to S^1$ given by $p(z) = z^n$
    is a covering map.
\end{problem}

\begin{soln}
    Let $z = \cos\theta + i\sin\theta\in S^1$ for some real $\theta$. Then,
    \[
	p\inv(z) = \left\{
	    \cos\left(\frac{\theta+2\pi k}n\right)
	    + i\sin\left(\frac{\theta+2\pi k}n\right) : k=0,1,\dots,n-1
	\right\}.
    \]
    If $\eps = \frac{\pi}{100}$, then the neighborhood
    $U = \{\cos\alpha + i\sin\alpha : \abs{\theta-\alpha} < \eps\}$
    has pre-image
    \begin{equation*}
        p\inv(U) =
	\bigcup_{k=0}^{n-1} \left\{\cos\alpha + i\sin\alpha :
	\abs{\frac\theta n - \alpha} < \frac\eps n\right\}.
    \end{equation*}
    Let the argument of the union be $U_k$;
    since $\eps$ is sufficiently small, the $U_k$ are pairwise disjoint.
    We show that $p|_{U_k}\colon U_k\to U$ is a homeomorphism.
    By construction, this restricted map is bijective and continuous,
    and its inverse is continuous because $p|_{U_k}$ is
    an open map --- in fact, the unrestricted $p\colon S^1\to S^1$
    is already an open map, since it sends
    $\{\cos\alpha + i\sin\alpha : \alpha\in (a,b)\}$ for some $a,b\in\RR$
    to the open set $\{\cos\alpha + i\sin\alpha : \alpha\in(na,nb)\}$,
    and sets of this kind generate the topology of $S^1$.

    The work above shows that $p$ is a local homeomorphism $S^1\to S^1$,
    and of course $p$ is continuous and surjective
    (the latter of which can be seen by e.g. noting that $p\inv(z)$,
    which we described above, is nonempty for all $z\in S^1$).
\end{soln}

\pagebreak

\begin{problem}
    [Munkres 54.5]
    Consider the covering map $p\colon\RR\times\RR\to S^1\times S^1$
    of Example 4 of \S53. Consider the path
    $f(t) = (\cos2\pi t,\sin2\pi t)\times(\cos4\pi t,\sin4\pi t)$
    in $S^1\times S^1$. Sketch what $f$ looks like
    when $S^1\times S^1$ is identified with the donut surface $D$.
    Find a lifting $\ol f$ of $f$ into $\RR\times\RR$,
    and sketch it.
\end{problem}

\pagebreak

\begin{problem}
    [Munkres 54.6]
    Consider the maps $g,h\colon S^1\to S^1$ given by
    $g(z) = z^n$ and $h(z) = 1/z^n$.
    Compute the induced homomorphisms $g_*$ and $h_*$
    of the infinite cyclic group $\pi_1(S^1,b_0)$ into itself.
\end{problem}

\begin{soln}
    Since $S^1$ is path connected, we can, without loss of generality,
    fix $b_0 = 1$ (since the change-of-basepoint homomorphism
    is an isomorphism).
    Let $\alpha\colon[0,1]\to S^1$ be an arbitrary loop in $S^1$ based at $1$
    and let $p\colon\RR\to S^1$ be the standard covering map.
    Suppose $\alpha$ lifts to a path $\widetilde \alpha$ in $\RR$
    starting at $0$; then, $g\circ\alpha$ lifts to
    the path $n\cdot\widetilde\alpha$ because, for all $t\in[0,1]$,
    \begin{equation*}
        p(n\cdot\widetilde \alpha(t)) = \cos(n\cdot\widetilde\alpha(t))
	+ i\sin(n\cdot\widetilde\alpha(t))
	= (\cos(\widetilde\alpha(t)) + i\sin(\widetilde\alpha(t)))^n
	= \alpha(t)^n = g(\alpha(t)),
    \end{equation*}
    where the second-to-last equality is by definition of
    $\widetilde\alpha$ being a lift of $\alpha$.
    Letting $\phi\colon\pi_1(S^1,b_0)\to\ZZ$ be the lifting correspondence
    used in Munkres's Theorem 54.5 (which is an isomorphism), we see that
    \begin{equation*}
	\phi([g\circ\alpha]) = n\phi([\alpha]) = \phi(n[\alpha]),
    \end{equation*}
    pulling $n$ into $\phi$ because $\phi$ is a homomorphism;
    and since $\phi$ is an isomorphism, applying its inverse
    shows that $\boxed{g_*([\alpha]) = [g\circ\alpha] = n[\alpha]}$.

    Similarly, $h(z) = 1/z^n = z^{-n}$ is pretty much the same map ---
    using $-n$ in place of $n$ above shows that $\boxed{h_*([\alpha]) = -n[\alpha]}$
    for all $[\alpha]\in\pi_1(S^1,b_0)$.
\end{soln}

\begin{problem}
    Generalize the proof of Theorem 54.5 to show that
    the fundamental group of the torus is isomorphic to the group $\ZZ\times\ZZ$.
\end{problem}

\begin{soln}
    Let $p\colon\RR\times\RR\to S^1\times S^1$ be the covering map
    in Munkres's \S53, Example 4, let $e_0 = (0,0)$,
    and let $b_0 = p(e_0)$. Then $p\inv(b_0)$ is the set
    $\ZZ\times\ZZ$ of ordered pairs of integers. Since $\RR^2$
    is simply connected, the lifting correspondence
    \begin{equation*}
        \phi\colon\pi_1(S^1\times S^1, b_0)\to\ZZ\times\ZZ
    \end{equation*}
    is bijective. We show that $\phi$ is a homomorphism
    (where $\ZZ^2$ has the componentwise additive structure),
    and the theorem is proved.

    Given $[f]$ and $[g]$ in $\pi_1(S^1\times S^1,b_0)$,
    let $\widetilde f$ and $\widetilde g$ be their respective liftings
    to paths in $\RR^2$ beginning at $(0,0)$.
    Let $(a,b) = \widetilde f(1)$ and $(c,d) = \widetilde g(1)$;
    then $\phi([f]) = (a,b)$ and $\phi([g]) = (c,d)$, by definition.
    Let $\widetilde{\widetilde g}$ be the path
    \begin{equation*}
	\widetilde{\widetilde g}(s) = (a,b) + \widetilde g(s)
    \end{equation*}
    in $\RR^2$. Because $p((a,b) + x) = p(x)$ for all $x\in\RR$,
    the path $\widetilde{\widetilde g}$ is a lifting of $g$;
    it begins at $(a,b)$. Then the product $\widetilde f * \widetilde{\widetilde g}$
    is defined, and it is the lifting of $f*g$ that begins at $(0,0)$.
    The endpoint of this path is $\widetilde{\widetilde g} = (a,b)+(c,d)$, i.e.
    \begin{equation*}
	\phi([f][g]) = (a,b) + (c,d) = \phi([f]) + \phi([g]).\qedhere
    \end{equation*}
\end{soln}

\pagebreak
\begin{problem}
    [Munkres 55.4]
    Suppose you are given the fact that for each $n$,
    there is no retraction $r\colon B^{n+1}\to S^n$. Prove the following:
    \begin{enumerate}[(a)]
        \ii The identity map $i\colon S^n\to S^n$ is not nullhomotopic.
	\ii The inclusion map $j\colon S^n\to\RR^{n+1} - \{\mb 0\}$
	is not nullhomotopic.
	\ii Every nonvanishing vector field on $B^{n+1}$
	points directly outward at some point of $S^n$,
	and directly inward at some point of $S^n$.
	\ii Every continuous map $f\colon B^{n+1}\to B^{n+1}$
	has a fixed point.
	\ii Every $(n+1)\times(n+1)$ matrix with positive real entries
	has a positive eigenvalue.
	\ii If $h\colon S^n\to S^n$ is nullhomotopic, then
	$h$ has a fixed point and $h$ maps some point $x$
	to its antipode $-x$.
    \end{enumerate}
\end{problem}

\begin{soln}
    (a) Suppose $i$ is nullhomotopic, and let $H\colon S^n\times I\to S^n$
    be a homotopy from $i$ to a constant map.
    We prove a generalization of $(1)\implies(2)$ in Lemma 55.3.
    Let $\pi\colon S^n\times I\to B^{n+1}$ be the map $\pi(x,t) = (1-t)x$.
    This is continuous, closed, and surjective, so it is a quotient map,
    collapsing $S^1\times\{1\}$ to the point $\mb 0$,
    and is otherwise injective.
    Since $H$ is constant on $S^1\times\{1\}$,
    it induces an extension $B^{n+1}\to S^n$ of $i$ via the quotient map $\pi$;
    this extension is exactly a retraction because
    it must be constant on $S^n$, in accordance with $i$.
    We are given that no retraction $B^{n+1}\to S^n$ exists,
    which is a contradiction, so $i$ is not nullhomotopic.
    \\

    (b) Let $r\colon\RR^{n+1}-\{\mb0\}\to S^n$ be the retraction
    $r(v) = v/\abs v$; then $i = r\circ j$ is the identity map,
    so if $j$ is homotopic to some constant map $c\colon S^n\to\RR^{n+1}$,
    then $i$ is homotopic to the constant map $r\circ c$,
    which contracts part (a).
    \\

    (c) Associate to each $x\in B^{n+1}$ the vector $v(x)$.
    Suppose first that $v(x)$ does not point directly inward at $x$
    for any $x\in S^n$. Let $w$ be the restriction of $v$ to $S^n$.
    Noting that $B^{n+1}$ is the quotient of $S^n\times I$
    under the quotient map $\pi$ from part (a), we find that
    $v\circ\pi\colon S^n\times I\to\RR^n$ is
    a homotopy of the map $w$ to a constant map
    since $v\circ\pi$ is constant on the fiber of the point $\mb 0\in B^{n+1}$.
    That is, $w$ is nullhomotopic.

    On the other hand, $w$ is homotopic to the inclusion map
    $j\colon S^n\to\RR^{n+1}-\{\mb0\}$ via the homotopy
    $F(x,t) = tx + (1-t)w(x)$; to verify that $F(x,t)$ is nowhere zero,
    we note that $F(x,t) = 0$ is equivalent to $w(x)$ being
    a negative scalar multiple of $x$, which we have assumed doesn't hold.
    It follows that $j$ is nullhomotopic, which contradicts the result of (b).

    The analogous result for an outward-pointing vector
    follows from applying the work above to the vector field $(x,-v(x))$.
    \\

    (d) Suppose that $f(x)\ne x$ for every $x\in B^{n+1}$.
    Then, defining $v(x) = f(x) - x$ gives us a nonvanishing vector field
    $(x,v(x))$ on $B^{n+1}$. But the vector field
    cannot point directly outward at any point $x$ of $S^1$,
    as that would mean $f(x) - x = ax$ for some positive real $a$,
    so that $f(x) = (1 + a)x$ would lie outside $B^{n+1}$.
    We thus arrive at a contradiction.
    \\

    (e) Let $A$ be the subset of $S^n$ with all nonnegative coordinates;
    $A$ is easily seen to be homeomorphic to the ball $B^n$.
    Then, treating the given matrix as a linear transformation
    $T$ with respect to the standard basis on $\RR^n$,
    the map $x\mapsto T(x)/\abs{T(x)}$ has a fixed point $x_0$ by part (d),
    implying $T(x_0) = \abs{T(x_0)}x_0$, so that $\abs{T(x_0)}$
    is a positive eigenvalue for $T$.
    \\

    (f) As proven in part (a), since $h$ is nullhomotopic,
    it has an extension $\widetilde h\colon B^{n+1}\to S^n$.
    Since $\mb 0\notin S^n$, the vector field $(x,\widetilde h(x))$
    on $B^{n+1}$ is nonvanishing,
    so it points directly outward at some point of $S^n$
    and directly inward at some point of $S^n$.
    The former corresponds to a fixed point of $\widetilde h$
    and the latter to a point $x$ that gets mapped by $\widetilde h$ to $-x$
    because $\widetilde h$ takes on values in $S^n$ only.
\end{soln}










\end{document}

